% 제출본 전용 — 부록으로 붙는다. 양식의 '8. 부록(소스코드)'.
% 전체 소스는 제출 폴더 4_캡스톤디자인_실행파일소스 와 공개 저장소에 있다.

\section{주요 소스코드}

전체 소스코드는 제출 폴더와 공개 저장소에 있다\citep{repo}. 여기에는 본문의
주장을 직접 떠받치는 세 조각만 싣는다. 주석은 저장소의 것을 그대로 옮겼다.

\subsection{사건 시점 정렬}

2.4절에서 선행연구의 결론을 가른 것이 거래 날짜를 세는 법이라고 했다. 그 세는
법이 이 함수다. 달력 월이 아니라 \textbf{효력일에 맞춘 월}을 쓴다.

\begin{footnotesize}
\begin{verbatim}
def event_month(dates: pd.Series, event: str) -> pd.Series:
    """사건 효력일 기준 경과 개월. 0 = 효력일이 속한 '사건월'.

    달력 월이 아니라 **효력일에 맞춘 월**이다. 2020-06-23 지정이면
    한 달은 23일에 시작해 다음 달 22일에 끝난다.

        2020-06-22  ->  -1   (지정 전날)
        2020-06-23  ->   0   (지정 당일)
        2020-07-22  ->   0
        2020-07-23  ->   1

    달력 월로 자르면 6월 1~22일의 지정 전 거래가 k=0 에 섞여 들어간다.
    실제로 강남·송파 k=0 구간 3,343건 중 1,532건(45.8%)이 지정 전이었고,
    그만큼 처치 직후 효과가 희석되고 있었다 (백로그 6b).
    """
    e = EVENTS[event]["effect"]
    d = pd.to_datetime(dates) - pd.Timedelta(days=e.day - 1)
    return (d.dt.year - e.year) * 12 + (d.dt.month - e.month)
\end{verbatim}
\end{footnotesize}

\textbf{세 줄짜리 함수가 추정치를 39.6\,\%p 가른다.} 이 정렬을 쓰면 지정 직후
거래량 효과가 $-80.1\%$, 달력 월로 자르면 $-40.5\%$다. \textbf{그리고 이
함수는 거래마다 계약일이 있어야만 쓸 수 있다} --- 동 단위 월별 집계로는
애초에 불가능하다. 1층이 2층의 전제라는 말의 가장 짧은 증거다.
(\texttt{scripts/lib/policy.py})

\subsection{무작위화 추론}

처치 단위가 4개뿐이라 관행적 $p$ 값을 쓸 수 없다(5.3절). 대신 \textbf{처치
배정을 바꿔 가며 같은 추정을 전부 돌려} 진짜 배정이 몇 등인지 센다.

\begin{footnotesize}
\begin{verbatim}
def randomization(s: Spec, max_draws: int | None = None,
                  seed: int = SEED) -> tuple[np.ndarray, float, int]:
    """처치 동 수만큼을 고르는 모든 조합에 대해 보고할 효과를 낸다.

    조합이 `max_draws` 보다 많으면 **무작위로 그만큼만 뽑는다.** 강남·송파
    27개 동에서 4개를 고르면 17,550가지라 전수가 되지만, 서울 전역 177개
    동에서 5개를 고르면 10^10 을 넘어 전수가 불가능하다(백로그 21).
    진짜 배정은 반드시 포함시키므로 순위는 그대로 뜻이 있다.
    """
    nt = int(s.treated.sum())
    real = tuple(np.flatnonzero(s.treated))
    total = math.comb(s.G, nt)
    if max_draws is None or total <= max_draws:
        combos = list(itertools.combinations(range(s.G), nt))
    else:
        rng = np.random.default_rng(seed)
        picked = {real}
        while len(picked) < max_draws:
            picked.add(tuple(sorted(rng.choice(s.G, nt, replace=False).tolist())))
        combos = sorted(picked)
    yt = s.M(s.y)
    vals = np.empty(len(combos))
    for i, cmb in enumerate(combos):
        t_ = np.zeros(s.G, bool)
        t_[list(cmb)] = True
        Xt = s.M(s.make(t_)).T
        XX = Xt.T @ Xt
        # 가짜 동이 모두 한 가격군에만 단지를 가지면 삼중차분 계수가 식별되지
        # 않는다. 그런 배정은 비교에서 뺀다 (몇 개인지 보고한다).
        if np.linalg.matrix_rank(XX) < XX.shape[0]:
            vals[i] = np.nan
            continue
        vals[i] = s.effect(np.linalg.solve(XX, Xt.T @ yt))
    act = vals[combos.index(real)]
    return vals[~np.isnan(vals)], act, len(combos)
\end{verbatim}
\end{footnotesize}

주석의 마지막 두 줄이 이 코드에서 가장 조심한 부분이다. \textbf{가짜 배정 중에는
계수가 아예 식별되지 않는 것이 섞인다.} 그런 배정을 0으로 세면 진짜 배정의
순위가 부풀려진다. 빼고, 몇 개를 뺐는지 함께 보고한다.
(\texttt{scripts/lib/inference.py})

\subsection{산출물 신선도 점검}

3.7절의 유령 숫자 사건에서 만든 장치다. 전체 재실행은 \textbf{다시 만들어지는}
파일만 보므로, 아무도 만들지 않는 파일은 영원히 낡은 채로 남는다. 그래서
\textbf{바뀐 것이 아니라 안 바뀐 것}을 센다.

\begin{footnotesize}
\begin{verbatim}
# 파이프라인이 만드는 곳. 원자료와 공고 원문은 손으로 넣는 것이라 뺀다.
WATCH = [
    ("output/results", "*.json"),
    ("output/results", "*.csv"),
    ("output/figures", "*.png"),
    ("data/processed", "*"),
    ("data/policy", "*.csv"),
]
# 파이프라인이 만들지 **않는 것이 맞는** 파일. 앞부분이 맞으면 건너뛴다.
# 여기 적을 때는 왜 파이프라인 밖인지 한 줄로 남긴다.
SKIP = (
    ("data/policy/loan_events.csv", "손으로 정리한 표 (0015)"),
    ...
)
\end{verbatim}
\end{footnotesize}

\textbf{예외 목록에 이유를 적게 한 것이 요점이다.} 예외를 늘리는 것만으로 점검을
무력화할 수 있으므로, 왜 파이프라인 밖인지를 한 줄 쓰게 했다. 쓸 말이 없으면
그건 예외가 아니라 고아 파일이다. 붙이자마자 고아 파일 넷과 폐기된 값 둘을
잡았다. (\texttt{scripts/check\_stale.py})

\subsection{코드 규모}

\begin{table}[htbp]
  \centering
  \caption{소스코드 구성}
  \label{tab:v2-code}
  \begin{tabular}{lp{9.6cm}}
    \toprule
    \textbf{묶음} & \textbf{내용} \\
    \midrule
    \texttt{scripts/lib/} & 공용 --- 정책 이력, 추론, 경제 계산, 경로 처리,
      PDF 판독, JSON 경로 \\
    \texttt{scripts/} 수집·정제 & 실거래 수집, 정제, 공고 판독, 지정 구간 생성 \\
    \texttt{scripts/} 분석 & 사건연구 · 추론 · 이질성 · 파급 · 가격지수 ·
      강건성 \pipeAnalysis 단계 \\
    \texttt{scripts/replicate/} & 선행연구 재현 \pipeReplicate 편 \\
    \texttt{scripts/check\_*.py} & 자동 점검 넷 \\
    \texttt{run\_all.py} & 전체 진입점 --- \pipeSteps 단계 \\
    \bottomrule
  \end{tabular}
\end{table}
